% Transformative mathematical distillation prepared from a privately supplied
% transcription.  This file does not reproduce the source narrative.

\section{The Section 9 free-beam example: formalization specification}

The numerical example studies the two-dimensional zero-mode space of a
free-end fourth-derivative operator on the unit interval after adding the
bounded multiplication perturbation $u(t) \mapsto \varepsilon t u(t)$, with
$0<\varepsilon<100$.  The formal development is divided into a finite-moment
calculation and an analytic spectral certificate.  Only the latter requires a
construction of the closed differential operator and a proof that its third
eigenvalue exceeds $500$.

\subsection{Centered affine coordinates}

Put $x=2t-1$.  An affine function is represented by a pair $(a,b)$ denoting
$a+bx$.  The unit-interval moments needed by the example are
\[
 \langle(a,b),(c,d)\rangle=ac+\frac{bd}{3},
\]
\[
 \langle(a,b),t(c,d)\rangle
 =\frac{ac}{2}+\frac{ad+bc}{6}+\frac{bd}{6},
\]
and
\[
 \langle(a,b),t^2(c,d)\rangle
 =\frac{ac}{3}+\frac{ad+bc}{6}+\frac{2bd}{15}.
\]
The normalized zero modes correspond to
\[
 e_1=\left(\frac{\sqrt2}{2},-\frac{\sqrt2\sqrt3}{2}\right),
 \qquad
 e_2=\left(\frac{\sqrt2}{2}, \frac{\sqrt2\sqrt3}{2}\right).
\]
The first bilinear form proves that these vectors are orthonormal.  The second
gives the diagonal Ritz compression
\[
 \widehat A_0=\operatorname{diag}(\widehat\alpha_1,
 \widehat\alpha_2),\qquad
 \widehat\alpha_{1,2}
 =\frac{\varepsilon}{2}\left(1\mp\frac1{\sqrt3}\right).
\]
The third gives the initial residual Gram data
\[
 R^*R=\frac{\varepsilon^2}{30}
 \begin{pmatrix}11-\sqrt{75}&-1\\-1&11+\sqrt{75}\end{pmatrix}.
\]
Its characteristic roots are
\[
 \rho_{\pm}=\frac{\varepsilon^2}{30}(11\pm\sqrt{76}).
\]
After subtracting the Ritz compression, the orthogonal residual satisfies
\[
 \widehat R^*\widehat R=
 \frac{\varepsilon^2}{30}
 \begin{pmatrix}1&-1\\-1&1\end{pmatrix},
\]
so its only nonzero singular value is $\varepsilon/\sqrt{15}$ and each column
has norm $\varepsilon/\sqrt{30}$.

\subsection{Exact envelopes and printed decimals}

The formal statements retain exact radicals.  The decimal forms are rational
corollaries.  In particular,
\[
 \frac{1}{500}\sqrt{\frac{11+\sqrt{76}}{30}}<0.001622,
\]
\[
 \frac{1}{500}\left(
 \sqrt{\frac{11+\sqrt{76}}{30}}+
 \sqrt{\frac{11-\sqrt{76}}{30}}\right)<0.00218.
\]
These are the arithmetic components of the initial sine and Ky Fan two-norm
bounds.  The double-angle perturbation estimates use the exact coefficients
$2/500$ and $4/500$.

For the recentered residual, define
\[
 c_1=\frac12\left(1-\frac1{\sqrt3}\right),\qquad
 c_2=\frac12\left(1+\frac1{\sqrt3}\right).
\]
The exact tangent envelopes are
\[
 \frac{\varepsilon/(500\sqrt{15})}{1-(c_2/500)\varepsilon}
 \quad\hbox{and}\quad
 \frac{2\varepsilon/(500\sqrt{15})}{1-(c_2/500)\varepsilon}.
\]
Rational upper approximations produce the coefficients printed in the paper:
$0.0005164$, $0.0010328$, and $0.0015774$.
For a single Ritz vector the numerator improves to
$1/(500\sqrt{30})<0.0003652$, with denominator coefficient $c_1/500$ or
$c_2/500$.

\subsection{Historical lower-bound comparison}

The comparison method uses the two low roots of the symmetric arrowhead matrix
\[
 \begin{pmatrix}
 \widehat\alpha_1&0&\varepsilon/\sqrt{30}\\
 0&\widehat\alpha_2&\varepsilon/\sqrt{30}\\
 \varepsilon/\sqrt{30}&\varepsilon/\sqrt{30}&500
 \end{pmatrix}.
\]
Rather than formalizing an informal asymptotic remainder, the Lean interface
asks for certified roots of its exact characteristic polynomial.  Given a
sine-square estimate
\[
 s^2\leq\frac{\widehat\alpha-\check\alpha}
 {g-\check\alpha},
\]
a purely algebraic lemma yields
\[
 \frac{s^2}{1-s^2}\leq
 \frac{\widehat\alpha-\check\alpha}{g-\widehat\alpha}.
\]
This isolates the historical comparison from the Davis--Kahan perturbation
bounds.

\subsection{Domain limitation and explicit repair}

For the diagonal multiplier with entries $\mu^{-n}$, the geometric trial
sequence $\mu^n$ is mapped pointwise to the constant sequence.  Its first $N$
output coordinates therefore have squared energy exactly $N$, making the
domain failure visible through finite partial sums.  A concrete repair is to
truncate the trial sequence after $N$ coordinates.  The modified trial vector
has finite support, its image has finite support, and it agrees with the
original sequence on every prescribed finite prefix once $N$ is large enough.
This replaces the source's informal ``small modification'' by an explicit
family.

\subsection{Schur reduction and individual vectors}

For a block eigenproblem
\[
 \begin{pmatrix}A_0&B^*\\B&A_1\end{pmatrix}
 \binom{x}{y}=\lambda\binom{x}{y},
\]
any left inverse $C$ of $\lambda I-A_1$ gives
\[
 y=C(Bx),\qquad
 [A_0+B^*CB]x=\lambda x.
\]
This is formalized for arbitrary modules and linear maps.  In the numerical
example, the in-plane angle contribution has coefficient
$1/(2\sqrt{75})=\sqrt3/30$, while the complementary-coordinate contribution
has coefficient $1/\sqrt{15}=\sqrt{15}/15$.  Their Euclidean combination is
exactly
\[
 \sqrt{\left(\frac{\sqrt3}{30}\right)^2+
       \left(\frac{\sqrt{15}}{15}\right)^2}
 =\frac{\sqrt7}{10}.
\]
After division by the spectral denominator this gives the exact envelope whose
rational relaxation has numerator $0.00053$ and the two printed denominator
coefficients.

\subsection{Remaining analytic obligation}

The finite calculations do not prove that the free-end fourth derivative has a
self-adjoint closure with two-dimensional zero eigenspace and third eigenvalue
above $500$.  A complete analytic instantiation must construct that operator,
identify the two affine zero modes, establish positivity of multiplication by
$\varepsilon t$, transport the third-eigenvalue lower bound to the perturbed
operator, and connect the general Davis--Kahan theorem outputs to the exact
certificate fields.  The current Section 9 package makes every downstream
obligation explicit and proves the rest independently.
